\problem{1}
\textbf{[Отбор 2025, 7.1]} 
Натуральные числа $a$ и $b$ удовлетворяют равенству $a+b=\frac{1}{a}+\frac{1}{b}$. Найдите все значения, которые может принимать выражение $a^2+b^2$.

\problem{2}
\textbf{[Отбор 2024, 7.2]} 
При каком наибольшем натуральном $n$, где $n<100$, выражение

$$
\frac{\left(2^2-1\right)\left(3^2-1\right) \ldots\left(n^2-1\right)}{2 n}
$$


является квадратом натурального числа?

\problem{3}
\textbf{[Отбор 2024, 7.5]}
Найдите все целые значения, которые может принимать выражение

$$
\frac{p q+p^p+q^q}{p+q},
$$


где $p, q$ - простые числа.

\problem{4}
\textbf{[Отбор 2023, 7.2]}
Докажите, что для любых чисел $a, b, c, d$, сумма которых равна нулю, справедливо неравенство:

$$
\max (a, b)+\max (a, c)+\max (a, d)+\max (b, c)+\max (b, d)+\max (c, d) \geq 0 .
$$

\problem{5}
\textbf{[Отбор 2023, 7.4]}
Если в трехзначном числе $x$ переставить цифры в обратном порядке, то получится трехзначное число $y$. Известно, что произведение $x y$ оканчивается на два нуля. Найдите все возможные значения $x$.


\problem{6}
\textbf{[Отбор 2023, 7.5]}
Петя написал на 30 карточках некоторые натуральные числа так, что каждое из них не имеет простых делителей, больших 10 и положил их на стол в перевернутом виде. Вася должен взять наименьшее количество карточек так, чтобы среди выбранных всегда нашлись две, перемножив числа на которых получится квадрат натурального числа. Сколько карточек должен взять Вася?

\problem{7}
\textbf{[Отбор 2022, 7.2]}
На шести карточках написаны цифры $1,2,3,4,5,6$. Составьте из них шестизначное число $\overline{a b c d e f}$ такое, что $\overline{a b}$ делится на $2, \overline{a b c}$ делится на $3, \overline{a b c d}$ делится на $4, \overline{a b c d e}$ делится на $5, \overline{a b c d e f}$ делится на 6 . Сколько таких чисел можно получить?

\newpage

\hspace{3mm}

\problem{8}
\textbf{[Отбор 2022, 7.3]}
Во время подготовки к очному туру Миша усердно решал задачи из доступных ему курсов Сириуса. В некоторых задачах он мог себя проверить по ответу. В остальных задачах ответа не было. Число задач с ответом составляет $9 / 11$ от числа остальных. Миша решил $2 / 7$ задач с ответом и $3 / 5$ остальных. В итоге ему осталось решить больше 500 задач с ответом и меньше 500 остальных. Сколько задач было в доступных Мише курсах Сириуса?

\problem{9}
\textbf{[Отбор 2022, 7.4]}
Найдите все пары простых чисел ( $p, q$ ), для которых $15(p-1)(q-1)$ делится на $p q$.